% !TEX program = xelatex
\documentclass[11pt,a4paper]{ctexart}
\usepackage{amsmath,amssymb}
\usepackage[margin=1.65cm,top=1.5cm,bottom=1.6cm]{geometry}
\usepackage{enumitem}
\usepackage{tabularx}
\usepackage{xcolor}
\usepackage{fancyhdr}

\pagestyle{fancy}
\fancyhf{}
\renewcommand{\headrulewidth}{0pt}
\renewcommand{\footrulewidth}{0pt}
\fancyfoot[C]{\small\thepage}
\fancypagestyle{plain}{\fancyhf{}\renewcommand{\headrulewidth}{0pt}\fancyfoot[C]{\small\thepage}}

\setlength{\parindent}{0pt}
\setlength{\parskip}{0pt}
% --- 分页美观：禁止孤行、寡行；标题块整体不可分割 ---
\clubpenalty=10000
\widowpenalty=10000
\displaywidowpenalty=5000
\raggedbottom
\newif\ifjustheading
\justheadingfalse

\linespread{1.05}
\setlength{\abovedisplayskip}{3pt}
\setlength{\belowdisplayskip}{3pt}
\setlength{\abovedisplayshortskip}{2pt}
\setlength{\belowdisplayshortskip}{2pt}
\allowdisplaybreaks

% ---- 填空 ----
\newcommand{\bk}{\underline{\hspace{3.2em}}}

% ---- 选项块（悬挂左缩进，自动换行） ----
% ---- 选项：自适应分列（四列等宽对齐 -> 两列 -> 单列悬挂） ----
\newlength{\optA}\newlength{\optB}\newlength{\optC}\newlength{\optD}
\newlength{\optMax}\newlength{\optColW}
\newcommand{\optline}[1]{\par\addvspace{0.06em}\hangindent=1.4em\hangafter=1\noindent #1\par}
\newcommand{\opts}[4]{%
  \par\addvspace{0.1em}\begingroup
  \setlength{\tabcolsep}{0.5em}%
  \settowidth{\optA}{#1}\settowidth{\optB}{#2}\settowidth{\optC}{#3}\settowidth{\optD}{#4}%
  \setlength{\optMax}{\optA}%
  \ifdim\optB>\optMax\setlength{\optMax}{\optB}\fi
  \ifdim\optC>\optMax\setlength{\optMax}{\optC}\fi
  \ifdim\optD>\optMax\setlength{\optMax}{\optD}\fi
  \setlength{\optColW}{\linewidth}%
  \addtolength{\optColW}{-6\tabcolsep}\divide\optColW by 4 \addtolength{\optColW}{-0.35em}%
  \ifdim\optMax>\optColW
    \setlength{\optColW}{\linewidth}%
    \addtolength{\optColW}{-2\tabcolsep}\divide\optColW by 2 \addtolength{\optColW}{-0.35em}%
    \ifdim\optMax>\optColW
      \optline{#1}\optline{#2}\optline{#3}\optline{#4}%
    \else
      \begin{tabularx}{\linewidth}{@{}XX@{}}#1&#2\\#3&#4\end{tabularx}\par
    \fi
  \else
    \begin{tabularx}{\linewidth}{@{}XXXX@{}}#1&#2&#3&#4\end{tabularx}\par
  \fi
  \endgroup
}

% ---- 一~五级标题 ----
\newcommand{\Hone}[1]{\par\addvspace{0.4em}\begin{center}{\LARGE\bfseries #1}\end{center}\addvspace{0.1em}}
\newcommand{\Htwo}[1]{\par\ifjustheading\global\justheadingfalse\else\penalty-300\addvspace{0.9em plus 0.6em}\fi\noindent{\large\bfseries #1}\par\vspace{0.05em}\nobreak\hrule height 0.7pt\nobreak\vspace{0.3em}\nobreak\justheadingtrue}
\newcommand{\Hthree}[1]{\par\ifjustheading\global\justheadingfalse\else\penalty-300\addvspace{0.7em plus 0.6em}\fi\noindent{\normalsize\bfseries #1}\par\addvspace{0.3em plus 0.2em}\nobreak\justheadingtrue}
\newcommand{\Hfour}[2]{\par\ifjustheading\global\justheadingfalse\else\penalty-300\addvspace{0.55em plus 0.5em}\fi\noindent{\bfseries #1\enskip #2}\par\addvspace{0.25em plus 0.2em}\nobreak\justheadingtrue}
\newcommand{\Hfive}[2]{\par\ifjustheading\global\justheadingfalse\else\penalty-300\addvspace{0.4em plus 0.5em}\fi\noindent\hspace{0.7em}{\small\bfseries #1\enskip #2}\par\nobreak\justheadingtrue}

% ---- 题目：\Q{题号}{出处} / \Qn{题号} ----
\newcommand{\Q}[2]{\par\ifjustheading\global\justheadingfalse\nobreak\else\addvspace{0.42em plus 0.6em}\fi\noindent\textbf{【#1】}{\small［#2］}\hspace{0.25em}}
\newcommand{\Qn}[1]{\par\ifjustheading\global\justheadingfalse\nobreak\else\addvspace{0.42em plus 0.6em}\fi\noindent\textbf{【#1】}\hspace{0.25em}}

\begin{document}

\Hone{一、函数极限连续}
\begin{center}\small
编号：题号 \textbf{【1.x.y-n】} 中 $n$ 为本章题目连续序号、1.x.y 为该题所属三级条目；标题数字点分最多到三级，四级用大写字母（1.1.2-A），五级用括号数字（1.1.2-A-(1)）。\quad ［　］内为出处缩写；小题用 (1)(2) 分问。
\end{center}

\Htwo{1.1　函数、极限、连续}

\Hthree{1.1.1　函数性质}

\Hfour{1.1.1-A}{有界}

\Q{1.1.1-1}{04数三}函数 $f(x)=\frac{|x|\sin(x-2)}{x(x-1)(x-2)^{2}}$ 在下列哪个区间内有界（　）。
\opts{A. $(-1,0)$}{B. $(0,1)$}{C. $(1,2)$}{D. $(2,3)$}

\Hfour{1.1.1-B}{奇偶}

\Q{1.1.1-2}{02数二}设函数 $f(x)$ 连续，则下列函数中，必为偶函数的是（　）。
\opts{A. $\int_{0}^{x}f(t^{2})\mathrm{d}t$}{B. $\int_{0}^{x}f^{2}(t)\mathrm{d}t$}{C. $\int_{0}^{x}t[f(t)-f(-t)]\mathrm{d}t$}{D. $\int_{0}^{x}t[f(t)+f(-t)]\mathrm{d}t$}

\Q{1.1.1-3}{20数三}设奇函数 $f(x)$ 在 $(-\infty,+\infty)$ 上有一阶连续导数, 则（　）。
\opts{A. $\int_{0}^{x}[\cos f(t)+f'(t)]\mathrm{d}t$ 为奇函数}{B. $\int_{0}^{x}[\cos f(t)+f'(t)]\mathrm{d}t$ 为偶函数}{C. $\int_{0}^{x}[\cos f'(t)+f(t)]\mathrm{d}t$ 为奇函数}{D. $\int_{0}^{x}[\cos f'(t)+f(t)]\mathrm{d}t$ 为偶函数}

\Q{1.1.1-4}{24数一}设 $f(x) = \int_{0}^{x} \mathrm{e}^{\cos t} \mathrm{d}t, g(x) = \int_{0}^{\sin x} \mathrm{e}^{t^2} \mathrm{d}t$ ，则下列正确的是（　）。
\opts{A. $f(x)$ 为奇函数， $g(x)$ 为偶函数}{B. $f(x)$ 为偶函数， $g(x)$ 为奇函数}{C. $f(x), g(x)$ 均为奇函数}{D. $f(x), g(x)$ 均为周期函数}

\Hfour{1.1.1-C}{周期}

\Q{1.1.1-5}{08数一}设函数 $f(x)$ 是连续函数
(1) 利用定义证明函数 $F(x)=\int_{0}^{x}f(t)\mathrm{d}t$ 可导, 且 $F'(x)=f(x)$ ；
(2) 当 $f(x)$ 是以 2 为周期的周期函数时, 证明 $G(x) = 2\int_{0}^{x} f(t) \mathrm{d}t - x \int_{0}^{2} f(t) \mathrm{d}t$ 也是以 2 为周期的周期函数。

\Hfour{1.1.1-D}{综合}

\Q{1.1.1-6}{87数二}函数 $f(x)=|x\sin x|\mathrm{e}^{\cos x},x\in(-\infty,+\infty)$ ，是（　）。
\opts{A. 单调函数}{B. 周期函数}{C. 偶函数}{D. 有界函数}

\Hthree{1.1.2　极限}

\Hfour{1.1.2-A}{函数极限计算}

\Hfive{1.1.2-A-(1)}{0/0 型}

\Q{1.1.2-7}{92数一}求 $\lim_{x\to 0}\frac{\mathrm{e}^x - \sin x - 1}{1 - \sqrt{1 - x^2}}$。

\Q{1.1.2-8}{99数二}求 $\lim_{x\to 0}\frac{\sqrt{1 + \tan x} - \sqrt{1 + \sin x}}{x\ln(1 + x) - x^2}$。

\Q{1.1.2-9}{02数二}设 $y = y(x)$ 是二阶常系数微分方程 $y'' + py' + qy = e^{3x}$ 满足初始条件 $y(0) = y'(0) = 0$ 的特解，则当 $x \to 0$ 时，函数 $\frac{\ln(1 + x^2)}{y(x)}$ 的极限（　）。
\opts{A. 不存在}{B. 等于1}{C. 等于2}{D. 等于3}

\Q{1.1.2-10}{08数一/二}求极限 $\lim_{x\to0}\frac{[\sin x-\sin(\sin x)]\sin x}{x^{4}}$。

\Q{1.1.2-11}{08数三}求极限 $\lim_{x\to0}\frac{1}{x^{2}}\ln\frac{\sin x}{x}$。

\Q{1.1.2-12}{11数三}求极限 $\lim_{x\to0}\frac{\sqrt{1+2\sin x}-x-1}{x\ln(1+x)}$。

\Q{1.1.2-13}{09数三}$\lim_{x\to 0}\frac{\mathrm{e} - \mathrm{e}^{\cos x}}{\sqrt[3]{1 + x^2} - 1} =$ \bk。

\Q{1.1.2-14}{12数三}求极限 $\lim_{x\to 0}\frac{\mathrm{e}^{x^2}-\mathrm{e}^{2-2\cos x}}{x^4}$。

\Q{1.1.2-15}{14数二}设函数 $f(x)=\arctan x$ . 若 $f(x)=xf'(\xi)$ ，则 $\lim_{x\to0}\frac{\xi^{2}}{x^{2}}=(\bk)$。
\opts{A. 1}{B. $\frac{2}{3}$}{C. $\frac{1}{2}$}{D. $\frac{1}{3}$}

\Q{1.1.2-16}{15数一/三}$\lim_{x\to 0}\frac{\ln(\cos x)}{x^2} =$ \bk。

\Q{1.1.2-17}{25数一}$\lim_{x\to 0^{+}}\frac{x^x - 1}{\ln x\cdot\ln(1 - x)} =$ \bk。

\Hfive{1.1.2-A-(2)}{$\infty/\infty$ 型}

\Q{1.1.2-18}{91数二}$\lim_{x\to 0^{+}}\frac{1 - \mathrm{e}^{\frac{1}{x}}}{x + \mathrm{e}^{\frac{1}{x}}} =$ \bk。

\Q{1.1.2-19}{97数二}求极限 $\lim_{x\to -\infty}\frac{\sqrt{4x^2 + x - 1} + x + 1}{\sqrt{x^2 + \sin x}}$。

\Q{1.1.2-20}{07数三}$\lim_{x\to +\infty}\frac{x^3 + x^2 + 1}{2^x + x^3} (\sin x + \cos x) =$ \bk。

\Hfive{1.1.2-A-(3)}{$0\cdot\infty$ 型}

\Q{1.1.2-21}{96数二}$\lim_{x\to \infty}x\left\{\sin \left[\ln \left(1 + \frac{3}{x}\right)\right] - \sin \left[\ln \left(1 + \frac{1}{x}\right)\right]\right\} =$ \bk。

\Q{1.1.2-22}{18数二}$\lim_{x\to +\infty}x^{2}[\arctan (x + 1) - \arctan x] =$ \bk。

\Q{1.1.2-23}{23数三}$\lim_{x\to \infty}x^{2}\left(2 - x\sin \frac{1}{x} -\cos \frac{1}{x}\right) =$ \bk。

\Q{1.1.2-24}{26数三}$\lim_{x\to 0}\frac{1}{x}\left(\frac{\sqrt{1 + x^2}}{\sin x} -\frac{1}{\tan x}\right) =$ \bk。

\Hfive{1.1.2-A-(4)}{$\infty-\infty$ 型}

\Q{1.1.2-25}{90数三}极限 $\lim_{n\to \infty}\left(\sqrt{n + 3\sqrt{n}} -\sqrt{n - \sqrt{n}}\right) =$ \bk。

\Q{1.1.2-26}{20数一}$\lim_{x\to 0}\left[\frac{1}{e^x - 1} -\frac{1}{\ln(1 + x)}\right] =$ \bk。

\Hfive{1.1.2-A-(5)}{$0^{0}$ 型或 $\infty^{0}$ 型}

\Q{1.1.2-27}{88数二}$\lim_{x\to 0^{+}}\left(\frac{1}{\sqrt{x}}\right)^{\tan x} =$ \bk。

\Q{1.1.2-28}{10数三}求极限 $\lim_{x\to+\infty}\left(x^{\frac{1}{x}}-1\right)^{\frac{1}{\ln x}}$。

\Hfive{1.1.2-A-(6)}{$1^{\infty}$ 型}

\Q{1.1.2-29}{91数三}求极限 $\lim_{x\to0}\left(\frac{\mathrm{e}^{x}+\mathrm{e}^{2x}+\cdots+\mathrm{e}^{nx}}{n}\right)^{\frac{1}{x}}$ ，其中 n 是给定的自然数。

\Q{1.1.2-30}{02数三}设常数 $a \neq \frac{1}{2}$ , 则 $\lim_{n \to \infty} \ln \left[ \frac{n - 2na + 1}{n(1 - 2a)} \right]^n =$ \bk。

\Q{1.1.2-31}{10数一}极限 $\lim_{x\to\infty}\left[\frac{x^{2}}{(x-a)(x+b)}\right]^{x}=(\bk)$。
\opts{A. 1}{B. e}{C. $e^{a-b}$}{D. $e^{b-a}$}

\Q{1.1.2-32}{11数一}求极限 $\lim_{x\to0}\left[\frac{\ln(1+x)}{x}\right]^{\frac{1}{e^{x}-1}}$。

\Q{1.1.2-33}{12数三}$\lim_{x\to \frac{\pi}{4}}(\tan x)^{\frac{1}{\cos x - \sin x}} =$ \bk。

\Q{1.1.2-34}{16数二/三}求极限 $\lim_{x\to0}\left(\cos2x+2x\sin x\right)^{\frac{1}{x^{4}}}$。

\Q{1.1.2-35}{19数二}$\lim_{x\to 0}(x + 2^x)^{\frac{2}{x}} =$ \bk。

\Q{1.1.2-36}{19数三}$\lim_{n\to \infty}\left[\frac{1}{1\cdot 2} +\frac{1}{2\cdot 3} +\dots +\frac{1}{n(n + 1)}\right]^n =$ \bk。

\Q{1.1.2-37}{22数二/三}$\lim_{x\to 0}\left(\frac{1 + \mathrm{e}^{x}}{2}\right)^{\cot x} =$ \bk。

\Hfive{1.1.2-A-(7)}{含 $\mathrm{e}$ 或绝对值的极限}

\Q{1.1.2-38}{00数一}求 $\lim_{x\to0}\left(\frac{2+\mathrm{e}^{\frac{1}{x}}}{1+\mathrm{e}^{\frac{4}{x}}}+\frac{\sin x}{|x|}\right)$。

\Hfive{1.1.2-A-(8)}{利用导数定义求极限（除式）}

\Q{1.1.2-39}{94数三}已知 $f'(x_0) = -1$ ，则 $\lim_{x \to 0} \frac{x}{f(x_0 - 2x) - f(x_0 - x)} =$ \bk。

\Q{1.1.2-40}{11数二/三}设函数 $f(x)$ 在 x=0 处可导, 且 $f(0)=0$ , 则 $\lim_{x\to0}\frac{x^{2}f(x)-2f(x^{3})}{x^{3}}=(\bk)$。
\opts{A. $-2f'(0)$}{B. $-f'(0)$}{C. $f'(0)$}{D. 0}

\Hfive{1.1.2-A-(9)}{利用导数定义求极限（乘式）}

\Q{1.1.2-41}{13数一}设函数 $f(x)$ 由方程 $y - x = \mathrm{e}^{x(1 - y)}$ 确定，则 $\lim_{n\to \infty}n\left(f\left(\frac{1}{n}\right) - 1\right) =$ \bk。

\Q{1.1.2-42}{13数三}设曲线 $y = f(x)$ 与 $y = x^2 - x$ 在点(1,0)处有公共切线，则 $\lim_{n\to \infty}nf\left(\frac{n}{n + 2}\right) =$ \bk。

\Q{1.1.2-43}{24数二}已知 $\left\{\begin{aligned} x &= 1 + t^{3} \\ y &= e^{t^{2}} \end{aligned}\right.$ ，则 $\lim_{x \to +\infty} x\left[f\left(2 + \frac{2}{x}\right) - f(2)\right] = (\bk)$。
\opts{A. 2e}{B. $\frac{4}{3}e$}{C. $\frac{2}{3}e$}{D. $\frac{e}{3}$}

\Hfive{1.1.2-A-(10)}{含变限积分的极限（$x\to0$）}

\Q{1.1.2-44}{02数三}求极限 $\lim_{x\to0}\frac{\int_{0}^{x}\left[\int_{0}^{u^{2}}\arctan(1+t)\mathrm{d}t\right]du}{x(1-\cos x)}$。

\Q{1.1.2-45}{17数二/三}求极限 $\lim_{x\to0^{+}}\frac{\int_{0}^{x}\sqrt{x-t}e^{t}\mathrm{d}t}{\sqrt{x^{3}}}$。

\Q{1.1.2-46}{16数一}$\lim_{x\to 0}\frac{\int_{0}^{x}t\ln(1 + t\sin t)\mathrm{d}t}{1 - \cos x^2} =$ \bk。

\Q{1.1.2-47}{21数一/二}求极限 $\lim_{x\to0}\left(\frac{1+\int_{0}^{x}e^{t^{2}}\mathrm{d}t}{e^{x}-1}-\frac{1}{\sin x}\right)$。

\Hfive{1.1.2-A-(11)}{含变限积分的极限（$x\to\infty$）}

\Q{1.1.2-48}{00数二}设函数 $S(x) = \int_{0}^{x}|\cos t|\mathrm{d}t$。
(1) 当 $n$ 为正整数, 且 $n\pi \leqslant x < (n + 1)\pi$ 时, 证明: $2n \leqslant S(x) < 2(n + 1)$
(2) 求 $\lim_{x\to+\infty}\frac{S(x)}{x}$。

\Q{1.1.2-49}{14数一/二/三}求极限 $\lim_{x\to+\infty}\frac{\int_{1}^{x}\left[t^{2}\left(\mathrm{e}^{\frac{1}{t}}-1\right)-t\right]\mathrm{d}t}{x^{2}\ln\left(1+\frac{1}{x}\right)}$。

\Hfive{1.1.2-A-(12)}{其他}

\Q{1.1.2-50}{91数三}下列各式中正确的是（　）。
\opts{A. $\lim_{x\to0^{+}}\left(1+\frac{1}{x}\right)^{x}=1$}{B. $\lim_{x\to0^{+}}\left(1+\frac{1}{x}\right)^{x}=e$}{C. $\lim_{x\to\infty}\left(1-\frac{1}{x}\right)^{x}=-e$}{D. $\lim_{x\to\infty}\left(1+\frac{1}{x}\right)^{-x}=e$}

\Q{1.1.2-51}{00数三}设对任意的 $x$ ，总有 $\varphi(x) \leqslant f(x) \leqslant g(x)$ ，且 $\lim_{x \to \infty} [g(x) - \varphi(x)] = 0$ ，则 $\lim_{x \to \infty} f(x)$（　）。
\opts{A. 存在且等于零}{B. 存在但不一定为零}{C. 一定不存在}{D. 不一定存在}

\Q{1.1.2-52}{06数三}$\lim_{n\to \infty}\left(\frac{n + 1}{n}\right)^{(-1)^n} =$ \bk。

\Hfour{1.1.2-B}{数列极限}

\Hfive{1.1.2-B-(1)}{数列极限计算（含 $n$ 的积分）}

\Q{1.1.2-53}{03数二}设 $a_{n} = \frac{3}{2}\int_{0}^{\frac{n}{n + 1}}x^{n - 1}\sqrt{1 + x^{n}}\mathrm{d}x,$ 则 $\lim_{n\to \infty}na_n =$ \bk。

\Q{1.1.2-54}{09数二}$\lim_{n\to \infty}\int_0^1\mathrm{e}^{-x}\sin nx\mathrm{d}x =$ \bk。

\Hfive{1.1.2-B-(2)}{数列极限计算（定积分定义求 $n$ 项和）}

\Q{1.1.2-55}{98数一}求 $\lim_{n\to\infty}\left(\frac{\sin\frac{\pi}{n}}{n+1}+\frac{\sin\frac{2\pi}{n}}{n+\frac{1}{2}}+\cdots+\frac{\sin\frac{n\pi}{n}}{n+\frac{1}{n}}\right)$。

\Q{1.1.2-56}{12数二}求极限 \[ \lim_{n\to\infty} n\left(\frac{1}{1+n^{2}}+\frac{1}{2^{2}+n^{2}}+\cdots+\frac{1}{n^{2}+n^{2}}\right) = \bk \]。

\Q{1.1.2-57}{17数一/二/三}求 $\lim_{n\to \infty}\sum_{k = 1}^{n}\frac{k}{n^2}\ln \left(1 + \frac{k}{n}\right)$。

\Q{1.1.2-58}{21数一/二}设函数 $f(x)$ 在区间 [0,1] 上连续，则 $\int_{0}^{1} f(x) \, \mathrm{d}x = (\bk)$。
\opts{A. $\lim_{n\to\infty}\sum_{k=1}^{n}f\left(\frac{2k-1}{2n}\right)\frac{1}{2n}$}{B. $\lim_{n\to\infty}\sum_{k=1}^{n}f\left(\frac{2k-1}{2n}\right)\frac{1}{n}$}{C. $\lim_{n\to\infty}\sum_{k=1}^{2n}f\left(\frac{k-1}{2n}\right)\frac{1}{n}$}{D. $\lim_{n\to\infty}\sum_{k=1}^{2n}f\left(\frac{k}{2n}\right)\frac{2}{n}$}

\Q{1.1.2-59}{25数二}$\lim_{n\to\infty}\frac{1}{n^{2}}\left[\ln\frac{1}{n}+2\ln\frac{2}{n}+\cdots+(n-1)\ln\frac{n-1}{n}\right]=$ \bk。

\Hfive{1.1.2-B-(3)}{数列极限计算（二重积分定义求 $n$ 项和）}

\Q{1.1.2-60}{10数一}$\lim_{n\to\infty}\sum_{i=1}^{n}\sum_{j=1}^{n}\frac{n}{(n+i)(n^{2}+j^{2})}=(\bk)$。
\opts{A. $\int_{0}^{1}\mathrm{d}x\int_{0}^{x}\frac{1}{(1+x)(1+y^{2})}\mathrm{d}y$}{B. $\int_{0}^{1}\mathrm{d}x\int_{0}^{x}\frac{1}{(1+x)(1+y)}\mathrm{d}y$}{C. $\int_{0}^{1}\mathrm{d}x\int_{0}^{1}\frac{1}{(1+x)(1+y)}\mathrm{d}y$}{D. $\int_{0}^{1}\mathrm{d}x\int_{0}^{1}\frac{1}{(1+x)(1+y^{2})}\mathrm{d}y$}

\Q{1.1.2-61}{26数二}$f(x,y)$ 在区域 $D=\left\{(x,y)\mid0\leqslant x\leqslant1,0\leqslant y\leqslant1\right\}$ 上连续，且 $f(x,y)=f(y,x)$ ，则 $\iint_{D}f(x,y)\mathrm{d}x\mathrm{d}y=(\bk)$。
\opts{A. $2\lim_{n\to\infty}\sum_{i=1}^{n}\sum_{j=n+1-i}^{n}f\left(\frac{i}{n},\frac{j}{n}\right)\frac{1}{n^{2}}$}{B. $\frac{1}{2}\lim_{n\to\infty}\sum_{i=1}^{n}\sum_{j=i}^{n}f\left(\frac{i}{n},\frac{j}{n}\right)\frac{1}{n^{2}}$}{C. $2\lim_{n\to\infty}\sum_{i=1}^{2n}\sum_{j=1}^{2n+1-i}f\left(\frac{i}{2n},\frac{j}{2n}\right)\frac{1}{n^{2}}$}{D. $\frac{1}{2}\lim_{n\to\infty}\sum_{i=1}^{2n}\sum_{j=1}^{i}f\left(\frac{i}{2n},\frac{j}{2n}\right)\frac{1}{n^{2}}$}

\Hfive{1.1.2-B-(4)}{概念题}

\Q{1.1.2-62}{12数二}设 $a_{n} > 0 (n = 1, 2, \cdots)$ ， $S_{n} = a_{1} + a_{2} + \cdots + a_{n}$ ，则数列 $\{S_{n}\}$ 有界是数列 $\{a_{n}\}$ 收敛的（　）。
\opts{A. 充分必要条件}{B. 充分非必要条件}{C. 必要非充分条件}{D. 既非充分也非必要条件}

\Q{1.1.2-63}{14数三}$\lim_{n\to\infty}a_n=a$ ，且 $a\neq0$ ，则当n充分大时有（　）。
\opts{A. $|a_{n}| > \frac{|a|}{2}$}{B. $|a_{n}| < \frac{|a|}{2}$}{C. $a_{n} > a - \frac{1}{n}$}{D. $a_{n} < a + \frac{1}{n}$}

\Q{1.1.2-64}{15数三}设 $\{x_{n}\}$ 是数列. 下列命题中不正确的是（　）。
\opts{A. 若 $\lim_{n\to\infty}x_{n}=a$ ，则 $\lim_{n\to\infty}x_{2n}=\lim_{n\to\infty}x_{2n+1}=a$}{B. 若 $\lim_{n\to\infty}x_{2n}=\lim_{n\to\infty}x_{2n+1}=a$ ，则 $\lim_{n\to\infty}x_{n}=a$}{C. 若 $\lim_{n\to\infty}x_{n}=a$ ，则 $\lim_{n\to\infty}x_{3n}=\lim_{n\to\infty}x_{3n+1}=a$}{D. 若 $\lim_{n\to\infty}x_{3n}=\lim_{n\to\infty}x_{3n+1}=a$ ，则 $\lim_{n\to\infty}x_{n}=a$}

\Hfive{1.1.2-B-(5)}{收敛发散问题}

\Q{1.1.2-65}{07数一/二}设函数 $f(x)$ 在 $(0, +\infty)$ 内具有二阶导数，且 $f''(x) > 0$ ，令 $u_n = f(n) (n = 1, 2, \cdots)$ ，则下列结论正确的是（　）。
\opts{A. 若 $u_1 > u_2$ ，则 $\{u_n\}$ 必收敛}{B. 若 $u_1 > u_2$ ，则 $\{u_n\}$ 必发散}{C. 若 $u_1 < u_2$ ，则 $\{u_n\}$ 必收敛}{D. 若 $u_1 < u_2$ ，则 $\{u_n\}$ 必发散}

\Q{1.1.2-66}{08数一/二}设函数 $f(x)$ 在 $(-\infty, +\infty)$ 内单调有界， $\{x_{n}\}$ 为数列，下列命题正确的是（　）。
\opts{A. 若 $\{x_{n}\}$ 收敛，则 $\{f(x_{n})\}$ 收敛}{B. 若 $\{x_{n}\}$ 单调，则 $\{f(x_{n})\}$ 收敛}{C. 若 $\{f(x_{n})\}$ 收敛，则 $\{x_{n}\}$ 收敛}{D. 若 $\{f(x_{n})\}$ 单调，则 $\{x_{n}\}$ 收敛}

\Q{1.1.2-67}{17数二}设数列 $\{x_{n}\}$ 收敛，则（　）。
\opts{A. 当 $\lim_{n\to \infty}\sin x_n = 0$ 时， $\lim_{n\to \infty}x_n = 0$}{B. 当 $\lim_{n\to \infty}\left(x_n + \sqrt{|x_n|}\right) = 0$ 时， $\lim_{n\to \infty}x_n = 0$}{C. 当 $\lim_{n\to \infty}\left(x_n + x_n^2\right) = 0$ 时， $\lim_{n\to \infty}x_n = 0$}{D. 当 $\lim_{n\to \infty}\left(x_n + \sin x_n\right) = 0$ 时， $\lim_{n\to \infty}x_n = 0$}

\Q{1.1.2-68}{22数一/二}设数列 $\{x_{n}\}$ 满足 $-\frac{\pi}{2} \leqslant x_{n} \leqslant \frac{\pi}{2}$ ，则（　）。
\opts{A. 若 $\lim_{n\to \infty}\cos (\sin x_n)$ 存在，则 $\lim_{n\to \infty}x_n$ 存在}{B. 若 $\lim_{n\to \infty}\sin (\cos x_n)$ 存在, 则 $\lim_{n\to \infty}x_n$ 存在}{C. 若 $\lim_{n\to \infty}\cos (\sin x_n)$ 存在, 则 $\lim_{n\to \infty}\sin x_n$ 存在, 但 $\lim_{n\to \infty}x_n$ 不一定存在}{D. 若 $\lim_{n\to \infty}\sin (\cos x_n)$ 存在, 则 $\lim_{n\to \infty}\cos x_n$ 存在, 但 $\lim_{n\to \infty}x_n$ 不一定存在}

\Q{1.1.2-69}{24数二}已知数列 $\{a_{n}\} (a_{n} \neq 0)$ , 若 $\{a_{n}\}$ 发散, 则（　）。
\opts{A. $\left\{a_{n} + \frac{1}{a_{n}}\right\}$ 发散}{B. $\left\{a_{n} - \frac{1}{a_{n}}\right\}$ 发散}{C. $\left\{\mathrm{e}^{a_n} + \frac{1}{\mathrm{e}^{a_n}}\right\}$ 发散}{D. $\left\{\mathrm{e}^{a_n} - \frac{1}{\mathrm{e}^{a_n}}\right\}$ 发散}

\Hfive{1.1.2-B-(6)}{极限最值}

\Q{1.1.2-70}{22数三}已知 $a_{n}=\sqrt[n]{n}-\frac{(-1)^{n}}{n}(n=1,2,\cdots)$ ，则 $\{a_{n}\}$（　）。
\opts{A. 有最大值,有最小值}{B. 有最大值, 没有最小值}{C. 没有最大值, 有最小值}{D. 没有最大值, 没有最小值}

\Hfive{1.1.2-B-(7)}{无穷小比阶}

\Q{1.1.2-71}{23数二}已知 $\{x_{n}\}, \{y_{n}\}$ 满足: $x_{1}=y_{1}=\frac{1}{2}, x_{n+1}=\sin x_{n}, y_{n+1}=y_{n}^{2}(n=1,2,\cdots)$ ，则当 $n\to\infty$ 时，（　）。
\opts{A. $x_{n}$ 是 $y_{n}$ 的高阶无穷小}{B. $y_{n}$ 是 $x_{n}$ 的高阶无穷小}{C. $x_{n}$ 与 $y_{n}$ 是等价无穷小}{D. $x_{n}$ 与 $y_{n}$ 是同阶但不等价的无穷小}

\Hfive{1.1.2-B-(8)}{单调有界}

\Q{1.1.2-72}{96数一}设 $x_{1}=10, x_{n+1}=\sqrt{6+x_{n}}(n=1,2,\cdots)$ ，试证数列 $\{x_{n}\}$ 极限存在，并求此极限。

\Q{1.1.2-73}{99数二}设 $f(x)$ 是区间 $[0, +\infty)$ 上单调减少且非负的连续函数, $a_{n} = \sum_{k=1}^{n} f(k) - \int_{1}^{n} f(x) \, \mathrm{d}x (n=1,2,\cdots)$ , 证明数列 $\{a_{n}\}$ 的极限存在。

\Q{1.1.2-74}{02数二}设 $0 < x_{1} < 3, x_{n+1} = \sqrt{x_{n}(3 - x_{n})} (n = 1, 2, \cdots)$ ，证明数列 $\{x_{n}\}$ 的极限存在，并求此极限。

\Q{1.1.2-75}{06数一/二}设 $0 < x_{1} < \pi, x_{n+1} = \sin x_{n}$ 。
(1) 证明: $\lim_{n\to\infty}x_{n}$ 存在, 并求值；
(2) 求 $\lim_{n\to\infty}\left(\frac{x_{n+1}}{x_n}\right)^{\frac{1}{x_n^2}}$。

\Q{1.1.2-76}{11数一/二}(1) 证明对任意的正整数 $n$ , 都有 $\frac{1}{n+1} < \ln\left(1 + \frac{1}{n}\right) < \frac{1}{n}$ 成立；
(2) 设 $a_{n}=1+\frac{1}{2}+\cdots+\frac{1}{n}-\ln n(n=1,2,\cdots)$ ，证明数列 $\{a_{n}\}$ 收敛。

\Q{1.1.2-77}{13数二}设函数 $f(x)=\ln x+\frac{1}{x}$ 。
(1) 求 $f(x)$ 的最小值；
(2) 设数列 $\{x_{n}\}$ 满足 $\ln x_{n} + \frac{1}{x_{n+1}} < 1$ , 证明 $\lim_{n \to \infty} x_{n}$ 存在, 并求此极限。

\Q{1.1.2-78}{18数一/二/三}数列 $\{x_{n}\}$ 满足：$x_{1}>0$，$x_{n}\mathrm{e}^{x_{n+1}}=\mathrm{e}^{x_{n}}-1\ (n=1,2,\cdots)$。证明 $\{x_{n}\}$ 收敛，并求 $\lim\limits_{n\to\infty}x_{n}$。

\Hfour{1.1.2-C}{已知极限相关}

\Hfive{1.1.2-C-(1)}{已知极限存在确定参数}

\Q{1.1.2-79}{98数二}确定常数 a, b, c 的值, 使 $\lim_{x \to 0} \frac{ax - \sin x}{\int_{b}^{x} \frac{\ln(1 + t^{3})}{t} \mathrm{d}t} = c (c \neq 0)$。

\Q{1.1.2-80}{01数三}已知 $f(x)$ 在 $(-\infty,+\infty)$ 内可导，且 $\lim_{x\to\infty}f'(x)=e,\lim_{x\to\infty}\left(\frac{x+c}{x-c}\right)^x=\lim_{x\to\infty}[f(x)-f(x-1)]$ ，求c的值。

\Q{1.1.2-81}{10数三}若 $\lim_{x\to0}\left[\frac{1}{x}-\left(\frac{1}{x}-a\right)e^{x}\right]=1$ ，则 a 等于（　）。
\opts{A. 0}{B. 1}{C. 2}{D. 3}

\Q{1.1.2-82}{11数二}已知函数 $F(x) = \frac{\int_0^x\ln(1 + t^2)\mathrm{d}t}{x^\alpha}$ . 设 $\lim_{x\to +\infty}F(x) = \lim_{x\to 0^+}F(x) = 0,$ 试求 $\alpha$ 的取值范围。

\Q{1.1.2-83}{13数一}已知极限 $\lim_{x\to0}\frac{x-\arctan x}{x^{k}}=c$ ，其中 c, k 为常数，且 $c\neq0$ ，则（　）。
\opts{A. $k=2, c=-\frac{1}{2}$}{B. $k=2, c=\frac{1}{2}$}{C. $k=3, c=-\frac{1}{3}$}{D. $k=3, c=\frac{1}{3}$}

\Q{1.1.2-84}{18数一}若 $\lim_{x\to 0}\left(\frac{1 - \tan x}{1 + \tan x}\right)^{\frac{1}{\sin kx}} = \mathbf{e},$ 则 $k =$ \bk。

\Q{1.1.2-85}{18数三}已知实数 a, b 满足 $\lim_{x\to+\infty}\left[(ax+b)\mathrm{e}^{\frac{1}{x}}-x\right]=2$ ，求 a, b。

\Q{1.1.2-86}{24数一}$\lim_{x\to 0}\frac{(1 + ax^2)^{\sin x} - 1}{x^3} = 6$ ，则 $a =$ \bk。

\Q{1.1.2-87}{21数三}已知 $\lim_{x\to 0}\left[a\arctan{\frac{1}{x}} + (1 + |x|)^{\frac{1}{x}}\right]$ 存在，求 $a$ 的值。

\Hfive{1.1.2-C-(2)}{已知一极限，求另一极限}

\Q{1.1.2-88}{00数二}若 $\lim_{x\to0}\frac{\sin6x+xf(x)}{x^{3}}=0$ ，则 $\lim_{x\to0}\frac{6+f(x)}{x^{2}}$ 为（　）。
\opts{A. 0}{B. 6}{C. 36}{D. $\infty$}

\Q{1.1.2-89}{16数三}已知函数 $f(x)$ 满足 $\lim_{x\to0}\frac{\sqrt{1+f(x)\sin2x}-1}{e^{3x}-1}=2,$ 则 $\lim_{x\to0}f(x)=$ \bk。

\Q{1.1.2-90}{20数三}设 $\lim_{x\to a}\frac{f(x)-a}{x-a}=b$ ，则 $\lim_{x\to a}\frac{\sin f(x)-\sin a}{x-a}=(\bk)$。
\opts{A. $b\sin a$}{B. $b\cos a$}{C. $b\sin f(a)$}{D. $b\cos f(a)$}

\Hfive{1.1.2-C-(3)}{已知极限推其他}

\Q{1.1.2-91}{22数一}已知 $f(x)$ 满足 $\lim_{x\to1}\frac{f(x)}{\ln x}=1$ ，则（　）。
\opts{A. $f(1)=0$}{B. $\lim_{x\to1}f(x)=0$}{C. $f'(1)=1$}{D. $\lim_{x\to1}f'(x)=1$}

\Hfour{1.1.2-D}{以极限形式给出的函数}

\Q{1.1.2-92}{15数二}函数 $f(x)=\lim_{t\to0}\left(1+\frac{\sin t}{x}\right)^{\frac{x^{2}}{t}}$ 在 $(-0,+\infty)$ 内（　）。
\opts{A. 连续}{B. 有可去间断点}{C. 有跳跃间断点}{D. 有无穷间断点}

\Q{1.1.2-93}{24数三}设函数 $f(x)=\lim_{n\to\infty}\frac{1+x}{1+nx^{2n}}$ ，则 $f(x)$（　）。
\opts{A. 在 x=1, x=-1 处都连续}{B. 在 x=1 处连续，在 x=-1 处不连续}{C. 在 x=1, x=-1 处都不连续}{D. 在 x=1 处不连续，在 x=-1 处连续}

\Hfour{1.1.2-E}{无穷小与无穷大的阶数}

\Hfive{1.1.2-E-(1)}{无穷小（概念题）}

\Q{1.1.2-94}{13数三}当 $x \to 0$ 时, 用“ $o(x)$ ”表示比 $x$ 高阶的无穷小量, 则下列式子中错误的是（　）。
\opts{A. $x \cdot o(x^2) = o(x^3)$}{B. $o(x) \cdot o(x^2) = o(x^3)$}{C. $o(x^2) + o(x^2) = o(x^2)$}{D. $o(x) + o(x^2) = o(x^2)$}

\Q{1.1.2-95}{22数二/三}若当 $x\to0$ 时，$\alpha(x)$，$\beta(x)$ 是非零无穷小量，则以下的命题中
(1) 若 $\alpha(x)\sim\beta(x)$，则 $\alpha^{2}(x)\sim\beta^{2}(x)$；
(2) 若 $\alpha^{2}(x)\sim\beta^{2}(x)$，则 $\alpha(x)\sim\beta(x)$；
(3) 若 $\alpha(x)\sim\beta(x)$，则 $\alpha(x)-\beta(x)=o(\alpha(x))$；
(4) 若 $\alpha(x)-\beta(x)=o(\alpha(x))$，则 $\alpha(x)\sim\beta(x)$。
真命题的序号为（　）。
\opts{A. (1)(3)}{B. (1)(4)}{C. (1)(3)(4)}{D. (2)(3)(4)}

\Q{1.1.2-96}{25数二}设函数 $f(x)$ , $g(x)$ 在 x=0 某去心邻域内有定义且恒不为 0, 若 $x \to 0$ 时, $f(x)$ 是 $g(x)$ 的高阶无穷小, 则当 $x \to 0$ 时（　）。
\opts{A. $f(x)+g(x)=o(g(x))$}{B. $f(x)g(x)=o(f^{2}(x))$}{C. $f(x)=o\left(\mathrm{e}^{g(x)}-1\right)$}{D. $f(x)=o\left(g^{2}(x)\right)$}

\Hfive{1.1.2-E-(2)}{无穷小（比阶）}

\Q{1.1.2-97}{13数二}设 $\cos x - 1 = x \sin \alpha(x)$ ，其中 $|\alpha(x)| < \frac{\pi}{2}$ ，则当 $x \to 0$ 时， $\alpha(x)$ 是（　）。
\opts{A. 比 x 高阶的无穷小量}{B. 比 x 低阶的无穷小量}{C. 与 x 同阶但不等价的无穷小量}{D. 与 x 等价的无穷小量}

\Q{1.1.2-98}{16数二}设 $\alpha_{1}=x(\cos\sqrt{x}-1)$ , $\alpha_{2}=\sqrt{x}\ln(1+\sqrt[3]{x})$ , $\alpha_{3}=\sqrt[3]{x+1}-1$ . 当 $x\to0^{+}$ 时, 以上 3 个无穷小量按照从低阶到高阶的排序是（　）。
\opts{A. $\alpha_{1},\alpha_{2},\alpha_{3}$}{B. $\alpha_{2},\alpha_{3},\alpha_{1}$}{C. $\alpha_{2},\alpha_{1},\alpha_{3}$}{D. $\alpha_{3},\alpha_{2},\alpha_{1}$}

\Q{1.1.2-99}{20数一/二}当 $x\to 0^{+}$ 时，下列无穷小中最高阶的是（　）。
\opts{A. $\int_0^x (\mathrm{e}t^2 -1)\mathrm{d}t$}{B. $\int_0^x\ln (1 + \sqrt{t^3})\mathrm{d}t$}{C. $\int_0^{\sin x}\sin t^2 \mathrm{d}t$}{D. $\int_0^{1 - \cos x}\sqrt{\sin^3t} \mathrm{d}t$}

\Hfive{1.1.2-E-(3)}{无穷小（已知高阶确定参数）}

\Q{1.1.2-100}{06数二}试确定 A, B, C 的值, 使得 $\mathrm{e}^{x}(1 + Bx + Cx^{2}) = 1 + Ax + o(x^{3})$ , 其中 $o(x^{3})$ 是当 $x \to 0$ 时比 $x^{3}$ 高阶的无穷小。

\Q{1.1.2-101}{14数二}当 $x \to 0^{+}$ 时, 若 $\ln^{\alpha}(1 + 2x), (1 - \cos x)^{\frac{1}{\alpha}}$ 均是比 x 高阶的无穷小量, 则 $\alpha$ 的取值范围是（　）。
\opts{A. $(2, +\infty)$}{B. $(1, 2)$}{C. $\left(\frac{1}{2}, 1\right)$}{D. $\left(0, \frac{1}{2}\right)$}

\Hfive{1.1.2-E-(4)}{无穷小（已知同阶确定参数）}

\Q{1.1.2-102}{96数一}设 $f(x)$ 有连续导数， $f(0)=0,\quad f'(0)\neq0,\quad F(x)=\int_{0}^{x}(x^{2}-t^{2})f(t)\mathrm{d}t$ ，且当 $x\to0$ 时， $F'(x)$ 与 $x^{k}$ 是同阶无穷小，则 k 等于（　）。
\opts{A. 1}{B. 2}{C. 3}{D. 4}

\Q{1.1.2-103}{19数一/二/三}当 $x \to 0$ 时, 若 $x - \tan x$ 与 $x^{k}$ 是同阶无穷小, 则 $k = (\bk)$。
\opts{A. 1}{B. 2}{C. 3}{D. 4}

\Q{1.1.2-104}{24数三}当 $x \to 0$ 时, $\int_{0}^{x} \frac{(1 + t^2) \sin t^2}{1 + \cos^2 t} \mathrm{d}t$ 与 $x^k$ 是同阶无穷小, 则 $k =$ \bk。

\Hfive{1.1.2-E-(5)}{无穷小（已知等价确定参数）}

\Q{1.1.2-105}{13数二/三}当 $x \to 0$ 时, $1 - \cos x \cdot \cos 2x \cdot \cos 3x$ 与 $ax^{n}$ 为等价无穷小量, 求 n 与 a 的值。

\Q{1.1.2-106}{15数一/二/三}设函数 $f(x)=x+a\ln(1+x)+bx\sin x, g(x)=kx^{3}$ . 若 $f(x)$ 与 $g(x)$ 在 $x\to0$ 时是等价无穷小, 求 a,b,k 的值。

\Q{1.1.2-107}{20数三}已知 $a, b$ 为常数，若 $\left(1 + \frac{1}{n}\right)^n - \mathbf{e}$ 与 $\frac{b}{n^a}$ 在 $n \to \infty$ 时是等价无穷小，求 $a, b$。

\Q{1.1.2-108}{23数一/二}当 $x \to 0$ 时, 函数 $f(x) = ax + bx^2 + \ln(1 + x)$ 与 $g(x) = e^{x^2} - \cos x$ 是等价无穷小, 则 $ab =$ \bk。

\Q{1.1.2-109}{26数二}已知当 $x \to 0$ 时, $ax^2 + bx + \arcsin x$ 与 $\sqrt[3]{1 + x^2} - 1$ 是等价无穷小, 则（　）。
\opts{A. $a = \frac{1}{3}, b = -1$}{B. $a = \frac{1}{3}, b = 1$}{C. $a = \frac{2}{3}, b = -1$}{D. $a = \frac{2}{3}, b = 1$}

\Hfive{1.1.2-E-(6)}{无穷大（无界与无穷大）}

\Q{1.1.2-110}{93数二}当 $x \to 0$ 时, $\frac{1}{x^2} \sin \frac{1}{x}$ 是（　）。
\opts{A. 无穷大}{B. 无穷小}{C. 有界但非无穷小}{D. 无界但非无穷大}

\Hfive{1.1.2-E-(7)}{无穷大（比阶）}

\Q{1.1.2-111}{10数三}设 $f(x)=\ln^{10}x, g(x)=x, h(x)=\mathrm{e}^{\frac{x}{10}}$ ，则当 x 充分大时有（　）。
\opts{A. $g(x) < h(x) < f(x)$}{B. $h(x) < g(x) < f(x)$}{C. $f(x) < g(x) < h(x)$}{D. $g(x) < f(x) < h(x)$}

\Hthree{1.1.3　连续}

\Hfour{1.1.3-A}{已知连续确定参数}

\Q{1.1.3-112}{18数二}设函数 $f(x)=\left\{\begin{aligned}-1,&x<0\\ 1,&x\geqslant0\end{aligned}\right.$ , $g(x)=\left\{\begin{aligned}2-ax,&x\leqslant-1,\\ x,&-1<x<0,\\ x-b,&x\geqslant0.\end{aligned}\right.$ 若 $f(x)+g(x)$ 在 R 上连续，则（　）。
\opts{A. a=3,b=1}{B. a=3,b=2}{C. a=-3,b=1}{D. a=-3,b=2}

\Q{1.1.3-113}{17数一/二/三}若函数 $f(x)=\left\{\begin{aligned}\frac{1-\cos\sqrt{x}}{ax},&x>0,\\ b,&x\leqslant0\end{aligned}\right.$ 在 x=0 处连续，则（　）。
\opts{A. $ab=\frac{1}{2}$}{B. $ab=-\frac{1}{2}$}{C. ab=0}{D. ab=2}

\Hfour{1.1.3-B}{间断点}

\Q{1.1.3-114}{10数二}函数 $f(x)=\frac{x^{2}-x}{x^{2}-1}\sqrt{1+\frac{1}{x^{2}}}$ 的无穷间断点的个数为（　）。
\opts{A. 0}{B. 1}{C. 2}{D. 3}

\Q{1.1.3-115}{13数三}函数 $f(x)=\frac{|x|^{x}-1}{x(x+1)\ln|x|}$ 的可去间断点的个数为（　）。
\opts{A. 0}{B. 1}{C. 2}{D. 3}

\Q{1.1.3-116}{20数二/三}函数 $f(x) = \frac{\mathrm{e}^{\frac{1}{x - 1}}\ln|1 + x|}{(\mathrm{e}^x - 1)(x - 2)}$ 的第二类间断点的个数为（　）。
\opts{A. 1}{B. 2}{C. 3}{D. 4}

\Q{1.1.3-117}{24数二}函数 $f(x)=|x|^{\frac{1}{(1-x)(x-2)}}$ 的第一类间断点的个数为（　）。
\opts{A. 3}{B. 2}{C. 1}{D. 0}

\end{document}